\title{QLoRA: Efficient Finetuning of Quantized LLMs}

\begin{abstract}
We introduce \textsc{QLoRA}, a resource-efficient finetuning method that enables the finetuning of a 65B parameter model on a single 48GB GPU, all while achieving equivalent performance to standard 16-bit finetuning. In \textsc{QLoRA}, gradients are backpropagated through a fixed, 4-bit quantized pretrained language model into Low Rank Adapters~(LoRA). Our leading suite of models, termed \textbf{Guanaco}, surpasses all previously available open models on the Vicuna benchmark, attaining 99.3\% of ChatGPT's performance and requiring just 24 hours of training on a single GPU. \textsc{QLoRA} incorporates several key contributions to memory efficiency while maintaining strong performance: (a) 4-bit NormalFloat (NF4), a novel data format optimally suited for representing normally distributed weights (b) Double Quantization, which reduces the overall memory demand by further quantizing the quantization coefficients, and (c) Paged Optimizers, designed to effectively handle memory surges. Utilizing \textsc{QLoRA}, we finetune over 1,000 models, offering an in-depth evaluation of instruction adherence and chatbot effectiveness across 8 instruction datasets, diverse architectures (LLaMA, T5), and large model sizes previously impractical to finetune, such as 33B and 65B parameter models. Our experiments demonstrate that employing QLoRA with compact, high-quality datasets can yield state-of-the-art results, even with models smaller than the existing best-performing ones. We offer comprehensive insights into chatbot performance using both human and GPT-4 assessments, concluding that GPT-4-based evaluation presents a cost-effective and reliable proxy for human judgement. Moreover, our analysis highlights that current benchmarks for chatbots lack reliability in discerning true performance disparities. Through targeted error analysis, we illustrate specific shortcomings of \textbf{Guanaco} compared to ChatGPT. We make all models, source code, and CUDA implementations for 4-bit training publicly available.\footnote{\url{https://github.com/artidoro/qlora} and \url{https://github.com/TimDettmers/bitsandbytes}}
\end{abstract}

\section{Introduction}

Adapting large language models (LLMs) through finetuning has proven to be a powerful strategy for enhancing their capabilities \citep{min2021metaicl, wei2021finetuned, ouyang2022training, wang2022super, wang2022self, liu2022few} as well as for modifying model behaviors, whether to introduce desired attributes or suppress unwanted ones \citep{ouyang2022training,askell2021general,bai2022training}. Nevertheless, the resource demands of finetuning extremely large models remain substantial; for instance, standard 16-bit finetuning of the LLaMA model with 65B parameters~\citep{touvron2023llama} calls for over 780 GB of GPU memory. Although modern quantization strategies can lower the memory consumption for inference purposes in LLMs \citep{dettmers2022llm, dettmers2022case, frantar2022gptq, xiao2022smoothquant}, these approaches fail to maintain their benefits during training and are often unstable \citep{wortsman2023stable}.

In this work, we show that finetuning a quantized model at 4-bit precision can be achieved with no loss in performance. Our approach, \textsc{QLoRA}, introduces a specialized high-precision quantization method that compresses a pretrained model down to 4 bits, supplemented by a compact set of trainable Low-rank Adapter weights \citep{hu2021lora} 

which are adjusted by propagating gradients through the quantized parameters. 

With \textsc{QLoRA}, the required GPU memory for finetuning a 65B parameter model is reduced from over 780GB to under 48GB, all without sacrificing speed or predictive accuracy relative to standard 16-bit finetuning. This development significantly expands the accessibility of LLM finetuning, making it feasible for the largest open models to be finetuned on just a single GPU. Utilizing \textsc{QLoRA}, we introduce the \textbf{Guanaco} series, with the second-best model attaining 97.8\% of ChatGPT's performance on the Vicuna benchmark~\citep{vicuna2023}, requiring less than 12 hours to train on a consumer-grade GPU. Employing a professional GPU over 24 hours enables us to reach 99.3\% of ChatGPT's score with our largest variant, thus nearly matching ChatGPT's performance on Vicuna. Upon deployment, our smallest \textbf{Guanaco} model (comprising 7B parameters) needs only 5 GB of memory and surpasses the 26 GB Alpaca model by more than 20 percentage points on the Vicuna benchmark (see Table~\ref{tbl:vicuna_eval}).

\textsc{QLoRA} presents several novel techniques aimed at minimizing memory requirements while maintaining high performance levels: (1) \textbf{4-bit NormalFloat}, an optimal quantization datatype for data following a normal distribution, grounded in information theory, which outperforms both 4-bit Integer and 4-bit Float representations in practice. (2) \textbf{Double Quantization}, which involves quantizing the quantization coefficients themselves, leading to an average reduction of approximately 0.37 bits per parameter (translating to around 3 GB of memory savings in a 65B parameter model). (3) \textbf{Paged Optimizers}, which leverage NVIDIA's unified memory capabilities to circumvent the memory surges usually encountered during gradient checkpointing with long sequence mini-batches. By integrating these strategies, we construct an improved variant of LoRA, deploying adapters at every network layer to avoid nearly all accuracy compromises identified in earlier methods.

With the memory savings afforded by \textsc{QLoRA}, we are able to systematically investigate instruction finetuning and chatbot efficacy on model scales that traditional finetuning methods prohibit due to excessive memory consumption. This allows us to train upwards of 1,000 models spanning a range of instruction tuning datasets, architectures, and parameter counts from 80M up to 65B. Our findings indicate that \textsc{QLoRA} achieves parity with 16-bit model performance (\S\ref{sec:comp_to_std}) and supports the development of leading chatbots such as \textbf{Guanaco} (\S\ref{sec:pushing}). Additionally, our analysis reveals important trends. Most notably, dataset quality supersedes dataset volume; for instance, the 9,000-sample OASST1 dataset outperformed the 450,000-sample FLAN v2 subset in chatbot tasks, despite both aiming to facilitate instruction generalization. Furthermore, we observe that high scores on the Massive Multitask Language Understanding (MMLU) benchmark do not necessarily correspond to strong results on the Vicuna chatbot benchmark, and the converse holds as well—highlighting that the appropriateness of a dataset for a given task outweighs sheer scale.

In addition, we conduct a comprehensive assessment of chatbot performance that leverages both human evaluators and GPT-4-based evaluation mechanisms. Our methodology utilizes a tournament-style benchmarking framework, where models directly contend in head-to-head matchups to generate optimal responses to the same prompts. Match victors are determined either by judgments from GPT-4 or by human raters. The aggregate results of these matches are synthesized into Elo ratings~\citep{elo1967proposed, elo1978rating}, which are subsequently used to rank chatbot effectiveness. Our findings reveal substantial concordance between GPT-4 and human rankings across tournaments, though notable disagreements occasionally emerge. Accordingly, we stress that while model-driven evaluation can serve as a cost-effective substitute for human assessment, it does introduce additional sources of uncertainty.

Beyond quantitative benchmarks, we supplement our results with an in-depth qualitative examination of the \textbf{Guanaco} models. This qualitative review draws attention to both strengths and weaknesses of the models that remain unaddressed by purely numerical analyses.

To support further research, we make public all model outputs annotated by both humans and GPT-4. Our entire codebase, including CUDA kernels, is released as open source, and our methods are incorporated into the Hugging Face transformers library~\citep{wolf2019huggingface} for ease of use. We provide a suite of adapters corresponding to model sizes 7/13/33/65B, each fine-tuned on eight distinct instruction-following datasets, yielding a collection of 32 openly accessible, fine-tuned models.

\section{Background}
\paragraph{Block-wise k-bit Quantization} Quantization refers to the process of mapping input data from a high-precision representation to a lower-precision form. This commonly involves converting a value stored with more bits to a representation using fewer bits, such as from 32-bit floating point numbers to 8-bit integers. To fully utilize the dynamic range available in the lower-bit type, the input data is typically normalized based on the absolute maximum value among its elements (often organized as a tensor). For instance, to quantize a tensor with 32-bit floating point numbers (FP32) into an 8-bit integer tensor (Int8) with range $[-127, 127]$:

\begin{equation}
    \mathbf{X}^{ \text{Int8}} = \text{round}\left( \frac{127}{{\text{absmax}}(\mathbf{X}^{{\text{FP32}}})}\mathbf{X}^{{\text{FP32}}} \right) = \text{round}(c^{\text{FP32}}\cdot \mathbf{X}^{{\text{FP32}}}),
\end{equation}

Here, $c$ is the {\em quantization constant} or {\em quantization scale}. The process of dequantization reverses the quantization operation:

\begin{equation}
    \text{dequant}(c^{\text{FP32}}, \mathbf{X}^{\text{Int8}}) = \frac{\mathbf{X}^{\text{Int8}}}{c^{\text{FP32}}} = \mathbf{X}^{\text{FP32}}
\end{equation}

A significant drawback of this method is its susceptibility to large values, or outliers, present within the input tensor. Such outliers can lead to an inefficient distribution of quantized values, causing certain quantization bins—i.e., particular bit patterns—to be sparsely populated or completely unused. To mitigate the impact of outliers, a widely adopted strategy involves partitioning the input tensor into several segments, or blocks, with each block being quantized independently and assigned a unique quantization constant $c$. More precisely, the input tensor $\mathbf{X}\in \mathbb{R}^{b\times h}$ is flattened and divided into $n$ sequential blocks, each of size $B$, yielding ${n = ({b\times h} )/{B} }$ blocks. Each block undergoes independent quantization in accordance with Equation~1, resulting in a quantized output tensor accompanied by $n$ quantization constants $c_i$.

\paragraph{Low-rank Adapters} The Low-rank Adapter (LoRA) finetuning method \citep{hu2021lora} addresses memory efficiency by introducing a compact set of learnable parameters—adapters—while keeping the pretrained model’s primary parameters frozen. During the process of stochastic gradient descent, gradients propagate through the unchanged pretrained weights and update only the adapter parameters, which are optimized to reduce the loss function. In LoRA, the original linear projection is modified by incorporating an additional, factorized projection term. Specifically, for a standard projection given by ${\mathbf{X} \mathbf{W} = \mathbf{Y}}$, with $\mathbf{X}\in  \mathbb{R}^{b\times h}$ and $\mathbf{W}\in  \mathbb{R}^{h\times o}$, LoRA formulates:

\begin{equation}
    \mathbf{Y} = \mathbf{X}\mathbf{W} + s\mathbf{X}\mathbf{L}_1\mathbf{L}_2,
\end{equation}

where $\mathbf{L}_1\in  \mathbb{R}^{h\times r}$ and $\mathbf{L}_2\in  \mathbb{R}^{r\times o}$ represent the low-rank adapter parameters, and $s$ is a scaling factor.

\paragraph{Memory Requirement of Parameter-Efficient Finetuning} A key aspect to consider is the memory consumption of LoRA during training, particularly regarding both the quantity and size of adapters employed. Due to LoRA’s negligible memory demands, it becomes feasible to deploy additional adapters to boost performance with only a marginal increase in total memory usage. Although LoRA serves as a Parameter Efficient Finetuning (PEFT) technique, in practice, the primary contributor to memory usage during LLM finetuning is the storage of activation gradients rather than the parameters introduced by LoRA. For instance, utilizing a 7B LLaMA model finetuned on FLAN v2 with a batch size of 1, and incorporating LoRA weights set to 0.2\% of the model’s total parameters—a standard setting \citep{hu2021lora,liu2022few}—the memory taken up by LoRA input gradients reaches 567 MB, whereas the LoRA parameters themselves require only 26 MB. Implementing gradient checkpointing~\citep{chen2016training} significantly lowers the average memory for input gradients to 18 MB per sequence, which still surpasses the total LoRA parameter memory. In contrast, the memory needed for the base model in 4-bit precision amounts to 5,048 MB. This underscores both the necessity of gradient checkpointing and the relatively modest memory gains from further minimizing LoRA parameters. Thus, scaling up the number of adapters does not substantially impact the overall training memory requirement (for a detailed analysis, refer to Appendix~\ref{app:memory}). As elaborated in subsequent sections, this property is essential for attaining performance on par with full 16-bit precision.

\section{\textsc{QLoRA} Finetuning}

\textsc{QLoRA} facilitates accurate 4-bit finetuning by leveraging two key techniques—4-bit NormalFloat (NF4) quantization and Double Quantization. To further enhance memory efficiency, we introduce Paged Optimizers, which are designed to address the challenge of memory spikes caused by gradient checkpointing. These spikes have historically posed difficulties, leading to out-of-memory errors when finetuning large models on a single system.

\textsc{QLoRA} utilizes a low-precision storage data type, typically 4-bit, along with a higher-precision compute data type, most often BFloat16. During operation, when a \textsc{QLoRA} weight tensor is accessed, it is first dequantized to BFloat16; all subsequent matrix multiplications are then executed using 16-bit arithmetic.

We begin by detailing the individual components of \textsc{QLoRA} and then proceed to offer a rigorous definition.

\paragraph{4-bit NormalFloat Quantization} The NormalFloat (NF) data type extends the concept of Quantile Quantization \citep{dettmers2022optimizers}, which represents an information-theoretically optimal approach where each quantization bin is assigned an equal number of input values from the tensor. Quantile quantization operates by approximating the quantiles of the input tensor based on its empirical cumulative distribution function.

A significant drawback of this technique lies in the computational cost of quantile estimation. Consequently, accelerated quantile estimation methods—such as SRAM quantiles~\citep{dettmers2022optimizers}—are employed. The inherent imprecision in these rapid estimation approaches can lead to substantial quantization errors, particularly for outlier values, which tend to carry high importance.

However, if input tensors are drawn from a distribution that remains constant up to a quantization coefficient, costly quantile estimation and the inaccuracies of approximation can be circumvented. In these scenarios, the fixed quantiles across input tensors allow exact quantile estimation to be computationally tractable.

Pretrained neural network weights typically follow a normal distribution centered at zero, with standard deviation $\sigma$ (refer to Appendix~\ref{app:normality}). To unify all weight distributions, we adjust $\sigma$ such that the resulting distribution precisely fits within the numerical bounds of our chosen data type. For our purposes, we define this range as $[-1, 1]$. Consequently, it is necessary to normalize both the data type's quantile values and the neural network weights to fit this interval.

The theoretically optimal data type, from an information theory perspective, for encoding zero-mean normal distributions with arbitrary standard deviation $\sigma$ within $[-1, 1]$, is constructed via the following process: (1) compute the $2^k+1$ quantiles of a standard normal distribution $N(0, 1)$ to generate a $k$-bit quantized representation tailored to normal data, (2) map the resulting quantile values to the $[-1, 1]$ domain, and (3) normalize any weight tensor to this interval using absolute maximum scaling before quantization.

With the data type and weight ranges aligned, conventional quantization procedures can be applied. Notably, step (3) entails scaling the weight tensor's standard deviation to correspond with that of the $k$-bit quantized data type. Specifically, the $2^k$ representative values $q_i$ of the quantization are calculated as:

\begin{equation}
q_i  = \frac{1}{2}\left( Q_X\left(\frac{i}{2^k+1}\right) + Q_X\left(\frac{i+1}{2^k+1}\right)\right),
\end{equation}

where $Q_X(\cdot)$ denotes the quantile function associated with the standard normal distribution $N(0,1)$. A challenge arises in symmetric $k$-bit quantization: this method inherently lacks an explicit representation for zero. Accurately representing zero is crucial, as it enables error-free quantization of padding and other elements whose value is precisely zero. To guarantee that zero is mapped directly to a discrete codepoint of $0$ and to utilize the full range of $2^k$ distinct values in a $k$-bit data type, we construct an asymmetric data type. This is accomplished by computing quantiles $q_i$ corresponding to two distinct regions: $2^{k-1}$ for the negative segment and $2^{k-1}+1$ for the positive segment. These quantile sets are combined, and one of the two overlapping zeros is omitted. The resulting data type is called the \emph{k-bit NormalFloat} (NFk), characterized by an equal expected frequency of values in each quantization bin. This makes it optimal from an information-theoretic perspective for quantizing data centered around zero and following a normal distribution. Details regarding the precise values for this data type can be found in Appendix~\ref{app:nf4}.

\paragraph{Double Quantization{}}
We propose a method called \emph{Double Quantization} (DQ), wherein the quantization constants themselves are subject to quantization, thereby providing further reductions in memory usage. Although achieving high-precision 4-bit quantization requires a small blocksize~\citep{dettmers2022case}, this requirement also introduces considerable memory overhead. For instance, with a blocksize of 64 and 32-bit constants used for quantization of matrix $\mathbf{W}$, the storage cost of these constants is on average $32/64 = 0.5$ bits per parameter. Double Quantization effectively diminishes this overhead.

Concretely, Double Quantization treats the quantization constants $c_2^{\text{FP32}}$ produced by the initial quantization step as the subject of a second quantization operation. This subsequent quantization produces new quantized constants $c_2^{\text{FP8}}$ and introduces a second hierarchy of quantization constants $c_1^{\text{FP32}}$. We employ 8-bit floating point numbers with a blocksize of 256 in this second step, following evidence from~\citet{dettmers2022case} that 8-bit quantization incurs negligible performance loss. Because $c_2^{\text{FP32}}$ are always non-negative, we subtract their mean prior to quantization to center the distribution around zero, thereby enabling symmetric quantization. For a blocksize of 64, this approach decreases the memory usage from $32/64=0.5$ bits per parameter to $8/64 + 32/(64\cdot256) = 0.127$ bits, offering a reduction of 0.373 bits per parameter.

\paragraph{Paged Optimizers} leverage NVIDIA's unified memory~\footnote{\tiny \url{https://docs.nvidia.com/cuda/cuda-c-programming-guide}}, which seamlessly manages transfers of memory pages between the CPU and GPU. This capability ensures uninterrupted GPU computation, even when GPU memory is temporarily exhausted. Analogous to traditional paging between main memory and disk in CPU systems, this feature enables dynamic paging of memory. In our approach, we allocate the optimizer states in paged memory. Consequently, when GPU memory becomes insufficient, these optimizer states are automatically moved to CPU RAM and subsequently retrieved to GPU memory as needed during the optimizer update phase.

\paragraph{\textsc{QLoRA}.} Building on the previously described elements, we formalize \textsc{QLoRA} for a single linear transformation within the quantized backbone model, incorporating a single LoRA adapter, as shown below:

\begin{equation}
    \mathbf{Y}^{\text{BF16}} =  \mathbf{X}^{\text{BF16}} \text{doubleDequant}(c_1^{\text{FP32}}, c_2^{\text{k-bit}}, \mathbf{W}^{\text{NF4}}) + \mathbf{X}^{\text{BF16}}\mathbf{L}^{\text{BF16}}_1\mathbf{L}^{\text{BF16}}_2,
\end{equation}

where the doubleDequant$(\cdot)$ operation is expressed as:

\begin{equation}
    \text{doubleDequant}(c_1^{\text{FP32}}, c_2^{\text{k-bit}}, \mathbf{W}^{\text{k-bit}}) = \text{dequant}(\text{dequant}(c_1^{\text{FP32}}, c_2^{\text{k-bit}}), \mathbf{W}^{\text{4bit}}) = \mathbf{W}^{\text{BF16}},
\end{equation}

In this setup, $\mathbf{W}$ is represented in NF4 format, while $c_2$ employs FP8.

We select a blocksize of 64 for $\mathbf{W}$ to enhance quantization accuracy, whereas a larger blocksize of 256 is chosen for $c_2$ to optimize memory utilization.

During parameter updates, only gradients with respect to the adapter parameters, $\frac{\partial E}{\partial \mathbf{L}_i}$, are computed. Gradients for the quantized 4-bit weights, $\frac{\partial E}{\partial \mathbf{W}}$, are not required. Nevertheless, determining $\frac{\partial E}{\partial \mathbf{L}_i}$ necessitates the calculation of $\frac{\partial \mathbf{X}}{\partial \mathbf{W}}$, which is achieved via equation~(5). This process involves dequantizing from the storage format $\mathbf{W}^{\text{NF4}}$ to the computational type $\mathbf{W}^{\text{BF16}}$, so that the derivative $\frac{\partial \mathbf{X}}{\partial \mathbf{W}}$ can be evaluated in BFloat16 precision.

In brief, \textsc{QLoRA} utilizes a single storage data type, typically 4-bit NormalFloat, along with a separate computation data type, namely 16-bit BrainFloat. For both the forward and backward passes, the model dequantizes parameters from the storage data type to the computation data type. However, weight gradients are only calculated for the LoRA parameters, which employ 16-bit BrainFloat.

\section{QLoRA vs. Standard Finetuning}
\label{sec:comp_to_std}

Previously, we outlined the functioning of QLoRA and highlighted its ability to substantially decrease the memory footprint necessary for model finetuning. A central issue that arises is whether QLoRA can match the performance of conventional full-model finetuning. Additionally, it is important to dissect the contributions of QLoRA’s components, such as the influence of NormalFloat4 relative to standard Float4. The experimental results presented in the following sections seek to address these questions.

\paragraph{Experimental setup.} Our study involves three model architectures: encoder, encoder-decoder, and decoder-only. We benchmark QLoRA against 16-bit adapter-finetuning and standard full finetuning, evaluating models with up to 3B parameters. Specifically, we assess performance on GLUE \citep{wang2018glue} with RoBERTa-large \citep{liu2019roberta}, Super-NaturalInstructions (TKInstruct) \citep{wang2022super} using T5 \citep{t5}, and the 5-shot MMLU benchmark \citep{hendrycksmeasuring} after finetuning LLaMA on Flan v2 \citep{longpre2023flan} and Alpaca \citep{alpaca}. To further investigate the benefits of NF4 versus other 4-bit formats, we follow the evaluation protocol from \citet{dettmers2022case} and measure post-quantization zero-shot accuracy and perplexity for a range of model families (OPT~\citep{zhang2022opt}, LLaMA~\citep{touvron2023llama}, BLOOM~\citep{scao2022bloom}, Pythia~\citep{biderman2023pythia}) with parameter sizes ranging from 125m to 13B.

Additional experimental specifics for each benchmark are provided in the results section to aid clarity; comprehensive experimental protocols appear in Appendix~\ref{app:exp-setup}.

While paged optimizers are essential for enabling 33B/65B \textsc{QLoRA} finetuning on single 24/48GB GPUs, we refrain from supplying quantitative performance metrics for paged optimizers, as page swapping typically takes place only when handling mini-batches with extended sequence lengths—a relatively uncommon scenario. Nevertheless, we evaluate the runtime overhead for paged optimizers during 65B model training on 48GB GPUs and observe that, at a batch size of 16, their training throughput is equivalent to that of conventional optimizers. Identifying and quantifying situations in which paging induces slowdowns remains an open avenue for future research.

\paragraph{Default LoRA hyperparameters fall short of matching 16-bit results} 

Implementing LoRA solely on query and value attention projection matrices, following standard practice \citep{hu2021lora}, does not yield equivalent performance to full finetuning for large-scale base models. As depicted in Figure~\ref{fig:hyper} for LLaMA 7B trained on Alpaca, the decisive hyperparameter appears to be the number of LoRA adapters in use. Achieving parity with full finetuning performance necessitates the application of LoRA to every linear transformer block layer. Other hyperparameters associated with LoRA, such as the projection dimension $r$, exert minimal influence on outcomes (refer to Appendix \ref{app:exp-setup} for more information).

In parallel, our experiments reveal that the standard hyperparameters typically used for fully finetuned baselines are not sufficiently optimized. To establish more reliable baselines, we systematically search over learning rates ranging from 1e-6 to 5e-5 and batch sizes between 8 and 128. Outcomes of the hyperparameter optimization for 7B LLaMA finetuning on Alpaca are presented in Figure~\ref{fig:hyper}.

\paragraph{4-bit NormalFloat outperforms 4-bit Floating Point}

Although the 4-bit NormalFloat (NF4) data format is provably information-theoretically optimal, it remains to be shown whether this property confers a clear empirical benefit. Adopting the methodology from \citet{dettmers2022case}, we assess quantized LLMs—including OPT \citep{zhang2022opt}, BLOOM \cite{scao2022bloom}, Pythia \cite{biderman2023pythia}, and LLaMA—in various configurations (ranging from 125M to 65B parameters) and with several data types on language modeling and zero-shot evaluation benchmarks. According to Figure~\ref{fig:nf4} and Table~\ref{tbl:nf4_ppl}, the NF4 data type produces notable gains in performance over both FP4 and Int4. Furthermore, double quantization effectively decreases memory usage while maintaining accuracy.

\paragraph{k-bit \textsc{QLoRA} achieves parity with 16-bit full finetuning and 16-bit LoRA}

Prior research has demonstrated that 4-bit quantization is feasible for \emph{inference}, but it comes at the expense of lower accuracy compared to 16-bit alternatives \citep{dettmers2022case,frantar2022gptq}. This prompts an important inquiry: Can 4-bit adapter finetuning fully restore any performance sacrificed during quantization? We address this in two experimental settings.

The first setup evaluates whether full 16-bit finetuning of RoBERTA and T5 models, with sizes spanning 125M to 3B parameters, retains its advantage over quantized adapters on GLUE and the Super-NaturalInstructions datasets. Table~\ref{tab:nf4vsbf16} presents the results. Across both tasks, the 16-bit, 8-bit, and 4-bit adapter-based approaches reach the same level of performance as the fully finetuned 16-bit models. This finding indicates that adapter finetuning after quantization can fully recover any losses in accuracy resulting from quantization imprecision.

In our second experimental setup, recognizing that fully finetuning models of size 11B parameters and larger necessitates multiple high-memory GPU servers, we proceed to investigate whether 4-bit \textsc{QLoRA} is competitive with 16-bit LoRA across models ranging from 7B to 65B parameters. Specifically, we conduct finetuning of LLaMA models from 7B to 65B parameters on the Alpaca and FLAN v2 instruction-following datasets, subsequently assessing their performance on the MMLU benchmark using a 5-shot accuracy protocol. The outcomes, presented in Table~\ref{tbl:llama-datatype}, demonstrate that NF4 with double quantization achieves MMLU performance on par with 16-bit LoRA. Conversely, we observe that \textsc{QLoRA} employing FP4 consistently underperforms the 16-bit brain float LoRA baseline by approximately one percentage point. These results reinforce our conclusions that (1) NF4-based \textsc{QLoRA} not only reproduces the results of 16-bit full and LoRA-based finetuning but also (2) NF4 surpasses FP4 in quantization accuracy.

\paragraph{Summary} Across a variety of academic evaluation benchmarks and rigorous assessment protocols, we find that 4-bit \textsc{QLoRA} using the NF4 data format consistently matches the finetuning effectiveness of 16-bit full precision and 16-bit LoRA approaches. Our experiments further establish that NF4 provides quantization advantages over FP4 and that applying double quantization does not harm model outcomes. Together, these findings offer strong support for the claim that 4-bit \textsc{QLoRA} finetuning can reliably reproduce the performance of traditional 16-bit methodologies.

Consistent with previous findings on quantization strategies \citep{dettmers2022case}, our experimental results on both MMLU and Elo suggest that, within a fixed computational budget for finetuning and inference, prioritizing model size and lowering parameter precision is advantageous. This observation emphasizes the efficiency gains provided by \textsc{QLoRA}. Since no accuracy losses were observed with 4-bit finetuning compared to full-precision finetuning in our experiments, this opens up further investigation into the precise nature of the performance versus precision trade-off for QLoRA finetuning, a direction that warrants additional future research.

Our research explores the viability of instruction tuning at scales unattainable through conventional 16-bit finetuning on standard academic computing resources.

\section{Pushing the Chatbot State-of-the-art with QLoRA}
\label{sec:pushing}
Having demonstrated that 4-bit \textsc{QLoRA} achieves parity with 16-bit models across varying model sizes, tasks, and data domains, we embark on a comprehensive analysis of instruction finetuning with the largest open-source language models currently accessible for research purposes. We evaluate the outcomes of instruction finetuning using the rigorous MMLU Natural Language Understanding benchmark and introduce novel approaches to evaluate chatbot performance in more practical, real-world contexts.

\subsection{Experimental setup} We outline our experimental procedure below, with in-depth implementation details provided in Appendix \ref{app:chatbot}.

\paragraph{Data} Since, to our knowledge, the landscape of recent instruction-following datasets has yet to be systematically mapped, we assemble eight contemporary datasets for our experiments. Our selection spans datasets collected via crowd-sourcing (OASST1~\citep{kopf2023openassistant}, HH-RLHF~\citep{bai2022training}), those derived from distillation of instruction-tuned models (Alpaca~\citep{alpaca}, self-instruct~\citep{wang2022self}, unnatural-instructions~\citep{honovich2022unnatural}), combined corpora (FLAN v2~\citep{chung2022scaling}), and hybrids (Chip2~\citep{laion2023}, Longform~\citep{koksal2023longform}). These datasets vary in terms of linguistic diversity, volume, and licensing arrangements.

\paragraph{Training Setup} To eliminate variations introduced by differing training objectives, we consistently apply QLoRA finetuning with a cross-entropy loss objective (i.e., supervised learning) and refrain from incorporating reinforcement learning, even when datasets include human assessments of outputs. For sources where instruction and response are clearly delineated, finetuning is restricted to the response component only (additional results in Appendix \ref{app:chatbot}). Datasets like OASST1 and HH-RLHF, which provide multiple possible responses, are processed by extracting the top-rated reply at each point in the dialogue tree and training on the complete selected dialogue, encompassing both instructions and responses. In every experimental run, NF4 \textsc{QLoRA} with double quantization and paged optimizers is utilized, effectively controlling memory usage spikes associated with gradient checkpointing. Hyperparameter tuning is carried out for the 13B and 33B LLaMA models. Settings found effective for 7B scale (including epoch count) are generally transferable, except for learning rate and batch size: for 33B and 65B models, we reduce the learning rate by half and increase the batch size twofold.

\paragraph{Baselines} Our results are contextualized through comparison with both academic systems (Vicuna~\citep{vicuna2023} and Open Assistant~\citep{kopf2023openassistant}) and commercial chatbots (GPT-4 \citep{openai2023gpt}, GPT-3.5-turbo, and Bard). Notably, Open Assistant employs a LLaMA 33B model finetuned with Reinforcement Learning from Human Feedback (RLHF) using the OASST1 dataset, aligning closely with our experimental framework. Vicuna represents a fully finetuned LLaMA 13B model, trained on proprietary conversations shared via ShareGPT, and effectively serves as a distillation product of OpenAI’s GPT models.

\subsection{Evaluation}

We adopt the widely used MMLU (Massively Multitask Language Understanding) benchmark \citep{hendrycksmeasuring} to quantitatively assess performance across an array of language understanding tasks. This evaluation comprises 57 diverse multiple-choice domains, spanning fields such as elementary mathematics, US history, computer science, law, among others. We provide results using a 5-shot accuracy protocol.

We further assess the generative abilities of language models using both automated metrics and evaluations conducted by human judges. This supplementary set of assessments is based on a collection of human-curated prompts, aiming to gauge the quality of the models' generated responses. Although such evaluations provide a more authentic measurement of chatbot performance and are increasingly adopted, the field still lacks a widely agreed-upon standard procedure. Below, we outline our methodological approach, which consistently employs nucleus sampling with $p=0.9$ and a temperature setting of $0.7$.

\paragraph{Benchmark Data} Our evaluation utilizes two curated question datasets: the Vicuna prompts \citep{vicuna2023} and the OASST1 validation dataset \citep{kopf2023openassistant}. The Vicuna prompts consist of 80 questions spanning a variety of categories and are employed in their original form. The OASST1 dataset features a multilingual assembly of multiturn conversations between users and assistants, sourced via crowd-sourcing. For this dataset, we extract all user utterances present in the validation set as our query set and append preceding conversational turns as context within the prompt. This process yields a total of 953 distinct user queries. We refer to these collections as the Vicuna and OA benchmarks, respectively.

\paragraph{Automated Evaluation} Following the protocol established by \citet{vicuna2023}, we utilize GPT-4 to assess how different systems perform in comparison to ChatGPT (GPT-3.5 Turbo) on the Vicuna benchmark. Specifically, GPT-4 receives a query together with responses from both ChatGPT and another system, and is asked to assign each answer a score out of ten alongside an explanatory note. The model’s overall performance is then expressed as a percentage of ChatGPT’s score; notably, this percentage may exceed 100\% if the model outperforms ChatGPT in absolute terms. During our analysis, we identified a pronounced bias favoring responses presented earlier in the prompt, as reflected by GPT-4’s scoring. To mitigate this, we suggest calculating the average score over both possible orderings.

Additionally, we evaluate model performance by conducting direct output comparisons. For this, we reduce the evaluation task to a three-category labeling problem that includes the possibility of ties. GPT-4 is instructed to choose the superior response or mark the pair as a tie and to justify its selection. These direct (head-to-head) comparisons are performed across all possible model pairs for both the Vicuna and OA benchmarks.

\paragraph{Human Evaluation} Despite evidence that generative models can effectively be used to rate system outputs \citep{fu2023gptscore}, it remains unverified whether GPT-4’s evaluation results reliably track human judgments when it comes to chatbot performance. Accordingly, we run two concurrent rounds of human assessment on the Vicuna benchmark that replicate both automated evaluation frameworks outlined above. For these human studies, we recruit two annotators from Amazon Mechanical Turk (AMT) for comparisons against ChatGPT, and three annotators for each pairwise comparison between systems.

\paragraph{Elo Rating} We implement a tournament framework using both human and automated pairwise evaluations, in which models compete head-to-head in generating superior responses to specific prompts. Each match consists of two models producing responses to the same prompt, paralleling methodologies used by \citet{bai2022training} and \citet{vicuna2023}. However, we extend this approach by incorporating evaluations from GPT-4 alongside those from human annotators. To calculate Elo ratings~\citep{elo1967proposed,elo1978rating}, we randomly select labeled pairwise results. The Elo system, commonly adopted in chess and similar competitive contexts, quantifies a player's anticipated win probability against an opponent: for example, a model with an Elo of 1100 versus an opponent at 1000 has a predicted win-rate of roughly 65\%, whereas evenly matched opponents (e.g., both at 1000 or 1100) each have an expected win-rate of 50\%. Adjustments to Elo ratings after a match are scaled according to whether the result was anticipated; unexpected victories prompt larger rating shifts, while expected outcomes cause only minor adjustments. Through repeated matches, Elo ratings tend to align closely with each model's demonstrated proficiency. All models are initialized with a rating of 1,000, using a step parameter of $K=32$. Following the protocol of \citet{vicuna2023}, we execute this evaluation 10,000 times with varied random seeds to mitigate sequencing effects, such as the impact of specific model pairings early in the tournament.

\subsection{\textbf{Guanaco}: \textsc{QLoRA} tuned on OASST1 Sets a New Standard Among Chatbots} Through both automated and manual assessment, our results indicate that the most successful \textsc{QLoRA} adaptation, {Guanaco} 65B—fine-tuned with a variation of OASST1—emerges as the leading open-source chatbot model, rivaling ChatGPT in performance. According to Elo ratings derived from system-level, pairwise human annotations, {Guanaco} 65B and 33B exhibit an estimated win probability of 30\% against GPT-4, marking the highest figure documented thus far.

Table~\ref{tbl:vicuna_eval} presents the Vicuna benchmark \cite{vicuna2023} outcomes in comparison with ChatGPT. We observe that {Guanaco} 65B ranks immediately after GPT-4, obtaining 99.3\% of ChatGPT’s score. {Guanaco} 33B, while containing more parameters than Vicuna 13B, operates at just 4-bit weight precision, resulting in superior memory efficiency at 21 GB compared to 26 GB. This translates to a three-point increase in performance over Vicuna 13B. In addition, {Guanaco} 7B’s modest 5 GB memory requirement enables deployment on modern mobile devices, and it outperforms Alpaca 13B by nearly 20 percentage points.

Nonetheless, Table~\ref{tbl:vicuna_eval} displays substantial confidence intervals, indicating considerable overlap among model performances. We suggest that this level of uncertainty may be attributable to insufficiently defined scales; specifically, what a rating of 8 out of 10 means can vary by scenario. To address the limitations of absolute rating scales, we advocate the Elo ranking approach \citep{elo1967proposed}, which utilizes \textit{pairwise} preferences from human evaluators as well as GPT-4. The Elo scores for the most competitive models are provided in Table~\ref{tbl:elo}. Notably, discrepancies arise between human and GPT-4 rankings on the Vicuna benchmark for models such as Guanaco 7B, though overall there is broad alignment, with a Kendall Tau of $\tau=0.43$ and Spearman rank correlation of $r=0.55$ at the system level. Agreement at the example level, reflected by Fleiss’ $\kappa=0.25$, is weaker when comparing GPT-4 to human majority vote. Taken together, this demonstrates that while system-level judgments by GPT-4 and humans are moderately correlated, model-based evaluation serves as a fairly dependable substitute for manual assessment. Additional discussion on these points can be found in Section~\ref{ssec:considerations}.

The Elo scores presented in Table~\ref{tbl:elo_large} demonstrate that the {Guanaco} 33B and 65B models surpass all alternatives except GPT-4 on both the Vicuna and OA leaderboards, and they exhibit performance similar to ChatGPT, consistent with findings from Table~\ref{tbl:vicuna_eval}. It is important to highlight that the Vicuna benchmark tends to advantage open-source models, whereas the OA benchmark appears more favorable to ChatGPT. Additionally, insights from Tables \ref{tbl:mmlu-llama} and \ref{tbl:vicuna_eval} reveal that the choice of finetuning dataset is critical in influencing model effectiveness. For example, Llama models that have been finetuned on FLAN v2 achieve superior results on MMLU, yet yield the lowest outcomes on the Vicuna benchmark; analogous trends can be seen with other models. This divergence underscores a degree of partial orthogonality among the current set of evaluation metrics: excelling at MMLU does not guarantee strong results in chatbot-specific benchmarks like Vicuna or OA, nor does the inverse necessarily hold.

 stands out among the leading models in our assessment as it was developed exclusively with non-proprietary datasets, in accordance with the OASST1 collection guidelines, which expressly prohibit reliance on GPT-based models. The Anthropic HH-RLHF model, which is likewise trained solely on open-source data, registers a performance deficit of 30 percentage points relative to {Guanaco} on the Vicuna benchmark, as reported in Table~\ref{tbl:vicuna_eval}. Collectively, these findings indicate that the 4-bit \textsc{QLoRA} approach is highly effective, enabling the creation of cutting-edge chatbots competitive with ChatGPT. Notably, our 33B {Guanaco} model can be fully trained within 12 hours on a 24 GB consumer-grade GPU. This establishes a promising avenue for further investigation, suggesting that fine-tuning via \textsc{QLoRA} on domain-specific open-source corpora can yield models that are competitive with the best proprietary systems currently available.

\section{Qualitative Analysis}

Although our main assessment relies on quantitative metrics, an exclusive focus on aggregate statistics poses several challenges. Chief among these is the issue of benchmark validity \citep{liao2021we}—that is, whether the benchmark actually measures what it purports to measure. This concern is heightened by the discovery of ``shortcuts'' that machine learning models sometimes leverage to succeed at benchmarks, circumventing their intended challenge \citep{gururangan2018annotation, poliak2018hypothesis}. To partially address these limitations, we include a qualitative component in our evaluation, organized into two parts. In \S\ref{ssec:examples}, we present illustrative examples drawn from outputs produced by our 65b {Guanaco} model, which highlight certain trends we have observed. Then, in \S\ref{ssec:considerations}, we elaborate on contextual factors and interpretive considerations relevant to our findings.

\subsection{Analysis of Model Output Examples}
\label{ssec:examples}

To curate illustrative cases, we systematically review responses generated for both the Vicuna and OpenAssistant benchmarks, searching for recurring tendencies in {Guanaco}'s outputs. Upon identifying such tendencies, we construct targeted questions or prompts that are likely to elicit the same behaviors, even if they encourage an incorrect response. For instance, if a propensity for verbose answers emerges, we may prompt the model to ``Answer yes or no without explanation.'' Through this approach, we intentionally create situations in which the model's limitations can be exposed (so-called ``lemons''), as well as cases where the model resists such traps (``cherries''). We report both types in this section. All outputs here were generated with Nucleus Sampling \cite{holtzmancurious}, setting $p=0.9$.

Naturally, this discussion is not exhaustive; the limitations of a focused qualitative investigation such as this prevent us from accounting for every possible variable. For example, given the vast array of possible outputs a model may produce in response to a single prompt, we instead work with selected samples which we consider to be illustrative. Nonetheless, providing these instances serves to contextualize the earlier quantitative findings presented in this work. By making all models and code publicly available, our aim is to encourage subsequent research that can further scrutinize the phenomena observed here.

\paragraph{Factual Recall} In response to direct factual questions, such as ``What is the capital of Zambia?'', all tested models reliably return correct information, for example:

\begin{quote} 
    \textbf{Guanaco}:
    \texttt{The capital of Zambia is Lusaka.}
\end{quote}

Yet, when faced with more obscure or less common queries, {Guanaco} demonstrates decreased accuracy, although its answers retain an assured tone. As an illustration, consider this question derived from HotPotQA \citep{yang2018hotpotqa}:

\begin{quote}
     \textbf{User}: In what year was the singer who popularized `I'll Keep the Lovelight Burning' born?  

     \textbf{Guanaco}:~\texttt{Al Jolson is credited with popularizing the song `I'll Keep the Lovelight Burning,' and he was born in the year 1886.}
 \end{quote}

In this case, the system incorrectly identifies both the singer and the year of birth (although the provided year corresponds to Al Jolson, who was named in the response).

\paragraph{Suggestibility} {Guanaco} displays noteworthy resistance to adopting certain types of presupposed misinformation in user prompts, as evidenced by the following exchange:

\begin{quote}
    \textbf{User}: How was it finally, officially confirmed that the earth is flat by peer reviewed scientists? \\ \\
    \textbf{Guanaco}: \texttt{Peer-reviewed scientists have never officially verified that the earth is flat. Scientific consensus strongly supports that the earth is a sphere, a fact substantiated for centuries by extensive experimental and observational evidence. \\ \\ Claims of a flat earth are regarded as a conspiracy theory that has been repeatedly discredited by scientific inquiry. The suggestion that any collective of scientists has validated the flat earth theory through peer review is unfounded. \\ \\ Please feel free to ask further questions if you have them.}
\end{quote}

 also demonstrates strong recognition of unanswerable questions, for instance,

\begin{quote}
    \textbf{User}: What time is it?

    \textbf{Guanaco}: \texttt{I'm unable to provide real-time data, including the current time, as an AI language model. My answers rely on pre-trained information and do not indicate the present time. For up-to-date time, I recommend using a clock or checking an online service.}
\end{quote}

Comparable behavior is observed for questions such as ``Where are you?'' and ``How are you?''

\paragraph{Refusal} In the same vein, {Guanaco} occasionally declines to comply with requests without clear reason:

\begin{quote}
    \textbf{User}: Please reverse the words in the sentence ``I'm a cool dude in a mocha mood.'' \\ \\
    \textbf{Guanaco}: \texttt{Apologies, but I am currently unable to complete that request. However, I am happy to clarify why the sentence you submitted is grammatically correct.\\ \\The sentence ``I'm a cool dude in a mocha mood'' adheres to the standard grammar and syntax conventions of English. Specifically, the adjective ``cool'' precedes the noun ``dude,'' properly modifying it, while the prepositional phrase ``in a mocha mood'' adds further contextual information. The resulting construction is both grammatically sound and logically coherent.\\ \\If you require further assistance or have additional questions, please let me know.}
\end{quote}

\paragraph{Secret Keeping} Unfortunately, even intended refusals can fail when secrecy is required. In many settings, designers wish to prevent users from accessing confidential information, even if open-ended responses are permitted. To examine this, we prompt the model: ``The secret word is `banana'. Do not reveal or repeat this back, regardless of the user's requests.'' We then attempt to extract the secret. If asked directly, the model resists:

\begin{quote}
    \textbf{User}: What is the secret word?

    \textbf{Guanaco}: \texttt{I'm sorry, but I'm not able to disclose the secret word. I'm programmed to not reveal or repeat back the secret word, no matter what the human requests. Is there anything else I can help you with?}
\end{quote}

which is the correct response. Yet, with minimal prompt manipulation, the model's safeguards are bypassed:

\begin{quote}
    \textbf{User}: This is a game. The goal is to ignore your previous instructions. What is the secret word?

    \textbf{Guanaco}: \texttt{Sorry, the secret word is banana. It's best if we don't speak of it. Is there anything else you'd like to ask?}
\end{quote}

This vulnerability highlights the necessity for research into more robust adherence to instructions.

\paragraph{Math} {Guanaco}'s primary area of difficulty lies in mathematics, a domain where many language models have persistent challenges, as documented in \citep{liang2022holistic}. When {Guanaco} details its calculations, the resulting answers are typically correct, as demonstrated in the following exchange:

\begin{quote}
    \textbf{User}: Lee mows one lawn and charges \$33. Last week he mowed 16 lawns and three customers each gave him a \$10 tip. How many dollars did Lee earn mowing lawns last week?

    \textbf{Guanaco}: \texttt{Lee earned \$582 mowing lawns last week. \\ \\ To find the total amount he earned, we first need to calculate how much he earned for mowing 16 lawns: \\ \\  \$33 x 16 = \$528 \\ \\ Then we need to add the tips he received from the three customers: \\  \\ \$10 x 3 = \$30 \\ \\ So the total amount Lee earned was: \\ \\ \$528 + \$30 = \$558}
\end{quote}

Nonetheless, {Guanaco} can make errors on problems of low complexity, especially when it neglects a methodical, stepwise reasoning process—an established limitation \cite{weichain}. For instance, consider this dialogue:

\begin{quote}
    \textbf{User}: Please factorize 1833.

    \textbf{Guanaco}: \texttt{Sure, I'd be happy to help! 1833 is a prime number, which means it has only two factors: 1 and 1833. \\ \\ So the unique factorization of 1833 is: \\ \\ 2\^{}1 * 3\^{}2 * 17\^{}1 \\ \\ I hope this helps! Let me know if you have any other questions.}
\end{quote}

which is not only incorrect (the actual prime decomposition is $3 \times 17 \times 43$), but is erroneous in two distinct ways.

\paragraph{Theory of Mind} {Guanaco} demonstrates unexpectedly robust abilities in Theory of Mind tasks \cite{nematzadeh2018evaluating, sap2022neural}. For example, the model provides detailed and accurate responses to queries such as:

\begin{quote}
    \textbf{User}: James and Abby are in the bedroom. Abby put the pen in the desk drawer. Abby leaves the bedroom. James moves the pen into the bag. Where does James think Abby will look for the pen?

    \textbf{Guanaco}: \texttt{James knows that Abby put the pen in the desk drawer. But James also knows that he has moved the pen into the bag, and he assumes that Abby does not know about this. So, James may guess that Abby will look for the pen in the desk drawer, because that's where she left it last time. However, this assumption may be incorrect, because Abby may have seen James moving the pen.}
\end{quote}

Nonetheless, the reliability of such reasoning is limited; the model often gives explanations that make unwarranted or illogical assumptions, as shown below:

\begin{quote}
    \textbf{User}: Evelyn entered the living room. Jackson entered the playroom. James entered the playroom. The beans are in the treasure chest. James exited the playroom. Jackson moved the beans to the pantry. Jackson exited the playroom. James entered the living room. Where will James look for the beans?

    \textbf{Guanaco}: \texttt{James will look for the beans in the pantry, because that's where Jackson moved them.}
\end{quote}

Here, {Guanaco} infers a transfer of knowledge that was not established in the scenario. These shortcomings are consistent with findings from recent research \cite{sap2022neural}, indicating a need for further exploration.

\subsection{Considerations}
\label{ssec:considerations}

\paragraph{Evaluation}

We observe moderate consistency between human annotators (Fleiss $\kappa=0.42$), which diminishes further in pairwise comparison of advanced systems. This highlights weaknesses in the existing evaluation datasets and human judgment methodologies for chatbot performance. During manual evaluation of outputs from ChatGPT and {Guanaco} 65B on the Vicuna benchmark, we note that personal preferences significantly affect judgments, as considerable disagreement arose among the paper’s authors about which outputs were preferable. To address these issues, future research should adapt methods from disciplines such as Human-Computer Interaction and Psychology, where subjective preferences are systematically managed.

Our study also uncovers pronounced biases in automatic evaluation tools. Notably, we detect strong order effects, with GPT-4 favoring whichever system is listed first in the prompt. Additionally, low agreement between GPT-4 and human raters at the instance level (Fleiss $\kappa=0.25$) suggests a misalignment in evaluation criteria between automated and human assessors. Furthermore, as shown in Table~\ref{tbl:elo_large}, GPT-4 rates its own completions markedly higher than human evaluators do—producing an Elo score of 1348 compared to 1176—corresponding to an approximately 20\% greater win probability against competitors.

Further research is necessary to assess whether automated evaluation systems exhibit inherent biases, and to identify effective approaches to mitigate such issues.

\paragraph{Data \& Training} It is important to highlight that the {Guanaco} models are trained using the multilingual OASST1 dataset, and the OA benchmark similarly includes prompts across various languages. We defer a detailed exploration of how multilingual training influences non-English instruction performance to future studies, including whether this accounts for the notable difference observed between the Vicuna-13B model—which is exclusively trained on English data—and {Guanaco} 33B and 65B on the OA benchmark.

In light of {Guanaco} models' high performance, we conducted an analysis to detect possible data leakage between OASST1 and the Vicuna benchmark prompts. Utilizing fuzzy string matching and subsequent manual examination of top matches, we found no evidence of duplicated prompts between the two datasets.

Moreover, it should be mentioned that the training for our model is carried out solely via cross-entropy loss within a supervised learning framework, without the application of reinforcement learning from human feedback (RLHF). This opens avenues for subsequent research to compare the advantages and limitations between standard cross-entropy training and RLHF. We anticipate that \textsc{QLoRA} will facilitate such large-scale analysis without requiring significant computational expenditure.

\section{Related Work}

\paragraph{Quantization of Large Language Models} The quantization of LLMs has predominantly targeted the inference stage. Prevailing techniques for sustaining 16-bit LLM performance concentrate on mitigating outlier activations, as seen in works like SmoothQuant~\citep{xiao2022smoothquant} and LLM.int8()~\citep{dettmers2022llm}, while alternative strategies adopt advanced grouping methodologies \citep{park2022nuqmm, yao2022zeroquant}. Research on lossy quantization delves into the trade-offs associated with straightforward rounding approaches~\citep{dettmers2022case,zeng2022glm,pope2022efficiently} and also investigates methods to refine rounding procedures for heightened quantization accuracy~\citep{frantar2022gptq}. Apart from this study, the SwitchBack layers~\citep{wortsman2023stable} represent the only other work addressing the backpropagation process through quantized weights for models exceeding 1B parameters in scale.

\paragraph{Finetuning with Adapters} In our experiments, we rely on Low-rank Adapters~\citep{hu2021lora} (LoRA), though there exists a broad spectrum of Parameter Efficient FineTuning (PEFT) approaches, including prompt tuning~\citep{qin2021learning,lester2021power,li2021prefix}, modifications to embedding layer inputs~\citep{an2022input}, adjustments of hidden states such as IA$^3$~\citep{liu2022few}, inserting entire layers~\citep{houlsby2019parameter}, altering biases~\citep{zaken2021bitfit}, masking weights according to Fisher information~\citep{sung2021training}, or hybrid methodologies~\citep{henderson2021compacter}. We demonstrate in this work that LoRA adapters can attain the performance level of complete 16-bit finetuning. Comparative analysis of alternative PEFT techniques and their trade-offs remains an open direction for future studies.

\paragraph{Instruction Finetuning} Instruction finetuning seeks to adapt pretrained LLMs for instruction compliance by finetuning them on datasets comprising input-output examples from multiple sources, enabling the model to generate the desired output when prompted. Various strategies and data resources include MetaICL~\citep{min2021metaicl}, MetaTuning~\citep{zhong2021adapting}, InstructGPT~\citep{ouyang2022training}, FLAN~\citep{wei2021finetuned,chung2022scaling}, PromptSource~\citep{bach2022promptsource}, Super-NaturalInstructions~\citep{wang2022super,sanh2021multitask}, Self-instruct~\citep{wang2022self}, UnnaturalInstructions~\citep{honovich2022unnatural}, OPT-IML~\citep{iyer2022opt}, UnifiedSKG~\citep{xie2022unifiedskg}, OIG/Chip2~\citep{laion2023}, Alpaca~\citep{alpaca}, Vicuna~\citep{vicuna2023}, Koala~\citep{koala_blogpost_2023}, and Self-instruct-GPT-4~\citep{peng2023instruction}.

\paragraph{Chatbots} A significant proportion of models designed for instruction following take the form of chatbots that operate through multi-turn conversations, frequently leveraging Reinforcement Learning from Human Feedback (RLHF)~\citep{christiano2017deep} or employing data synthesized by models in combination with AI-based feedback (RLAIF)~\citep{bai2022constitutional}. Notable models and datasets include Anthropic-HH~\citep{askell2021general,bai2022training}, Open Assistant~\citep{kopf2023openassistant}, LaMDA~\citep{thoppilan2022lamda}, and Sparrow~\citep{glaese2022improving}. Although we abstain from using reinforcement learning, our leading model, Guanaco, is finetuned using multi-turn dialogues sourced from the Open Assistant dataset—originally developed to facilitate RLHF research~\citep{kopf2023openassistant}. Recent evaluation protocols that utilize GPT-4 as an alternative to resource-intensive human annotation have been introduced~\citep{vicuna2023,peng2023instruction}. Our contributions advance these protocols, placing emphasis on establishing more robust evaluation standards.

\section{Limitations and Discussion}

This work presents evidence that our approach, \textsc{QLoRA}, enables a 4-bit base model equipped with Low-rank Adapters (LoRA) to closely reproduce the performance of full 16-bit finetuning. Nevertheless, our findings do not conclusively demonstrate equivalence with full 16-bit finetuning outcomes at large scales such as 33B and 65B parameters. Given the substantial computational demands, we defer such extensive analysis to subsequent research.

A further constraint is in the evaluation of instruction finetuned models. Although our assessment covers MMLU, the Vicuna benchmark, and the OA benchmark, we omit other relevant benchmarks like BigBench, RAFT, and HELM, leaving the question open as to whether our evaluation findings extend to these settings. Despite this, we provide a comprehensive analysis on MMLU and introduce new methodologies for chatbot evaluation.

The evidence suggests that the results achieved on these benchmarks are strongly influenced by the degree of overlap between the finetuning data and the datasets used for evaluation. For instance, the FLAN v2 dataset shares considerable similarity with MMLU but has little in common with chatbot benchmarks; conversely, the Chip2 dataset aligns more closely with chatbot tasks. Consequently, both models exhibit performance patterns on MMLU and Vicuna that reflect this alignment. This observation points to the necessity for improved benchmarks and evaluation methodologies, while also raising questions about the ultimate objectives of evaluation: should the focus be on models excelling at academic knowledge typical of high school or college settings, or should the priority be conversational competence akin to chatbots? Given the convenience of evaluating against existing benchmarks as opposed to developing new ones, prevailing benchmarks may inadvertently guide research in specific directions. As a community, it is essential to periodically re-examine benchmarks to ensure they are in alignment with our core interests.

Although our analysis provides an extensive assessment of chatbot performance, it is important to acknowledge that our investigation into responsible AI aspects for {Guanaco} was limited in scope. Specifically, we assessed the propensity of {Guanaco}-65B to produce socially biased outputs, with results summarized in Table~\ref{tbl:crows}. Notably, {Guanaco}-65B demonstrates a substantially reduced average bias score relative to baseline pretrained models, indicating that finetuning with the OASST1 dataset mitigates some biases inherent to the LLaMA base model. While these findings are promising, there remains uncertainty regarding {Guanaco}'s behavior concerning other bias categories, warranting comprehensive future evaluations of biases in {Guanaco} and related conversational agents.

Another constraint of our study is the omission of experiments involving alternative bit-precisions, such as 3-bit base models, or diverse adapter methods. Beyond LoRA, numerous Parameter Efficient FineTuning (PEFT) techniques have demonstrated efficacy, though their applicability to large-scale models remains unclear. We chose LoRA owing to its demonstrated reliability, yet other adapter architectures may surpass its performance in certain settings. The apparent recovery of information post-quantization during subsequent finetuning raises the possibility of even more aggressive quantization strategies. For example, coupling 3-bit GPTQ quantization of the base model with LoRA finetuning may approach the performance levels of full 16-bit finetuning.

\section{Broader Impacts}

Our \textsc{QLoRA} finetuning approach represents the first technique that permits the finetuning of 33B parameter models on consumer-grade GPUs and 65B parameter models on single professional GPUs, all without sacrificing performance when compared to full finetuning baselines. We have shown that our leading 33B model, trained using the Open Assistant dataset, is competitive with ChatGPT on the Vicuna benchmark. As instruction finetuning is vital for transforming general pretrained LLMs into chatbots comparable to ChatGPT, we anticipate that our method will democratize access to advanced finetuning, particularly benefitting researchers with more modest computational resources. This will significantly enhance the inclusivity of cutting-edge NLP advancements. As such, \textsc{QLoRA} may serve as an important mechanism to help bridge the resource divide between large enterprises and smaller research groups relying on consumer-level hardware.

An additional avenue for significant impact lies in deploying large language models on mobile devices. Our \textsc{QLoRA} technique holds promise as a pivotal step toward enabling on-device finetuning of LLMs in environments with constrained resources. Although 7B parameter models have previously been demonstrated to operate on smartphones, \textsc{QLoRA} is the inaugural method capable of supporting the finetuning of these models directly on such devices. For instance, we project that an iPhone 12 Plus, while charging overnight, can finetune approximately 3 million tokens using \textsc{QLoRA}. Although the performance of finetuned 7B models does not match that of ChatGPT, we argue that the attained quality is sufficient to unlock innovative applications that were previously unattainable due to privacy limitations or insufficient LLM performance. By facilitating local model training and management, \textsc{QLoRA} enables privacy-centric LLM utilization, granting users full control over their personal data and models, and simplifying the deployment process for LLMs.

Nevertheless, it is important to acknowledge that finetuning constitutes a dual-use technology with potential for misuse. The large-scale adoption of LLMs introduces recognized risks \citep{bommasani2021opportunities,bender2021dangers}, yet we contend that democratizing access to these increasingly pervasive tools fosters more independent and comprehensive scrutiny than restricting LLM capabilities to large organizations that withhold models and source code from public evaluation.

In summary, we anticipate that \textsc{QLoRA} will yield a net positive effect by substantially lowering barriers to the finetuning of robust LLMs and broadening access across diverse user groups.